%% kapa: exam=1 round=19 year=2008 period=1 form=A subject=accounting subject_ko=회계학 seq=39 number=119
%% source: Q-Net 감정평가사 기출문제 (cst00309) artlSeq=1001962 file=extracted/1008/1001962/2008년 제19회 감정평가사 1차시험 기출문제/1교시.hwp
%% conversion: kordoc 4.24.51 -> tools/gen1.py (auto); review=proofread-sonnet; flags=none
\begin{question}{119}
다음 중 재무회계의 기본개념에 대한 설명으로 가장 옳은 것은?
\begin{choices}
  \item 대차대조표는 기업의 경영활동 중 투자활동과 재무활동의 결과뿐만 아니라 연중 발생한 영업활동의 내용까지 나타낸다.
  \item 재무제표는 발생기준에 따라 작성된다. 발생주의 회계는 재무회계의 기본적 특징으로서 재무제표의 기본요소의 정의 및 인식, 측정과 관련이 있다. 따라서 현금흐름표도 발생기준에 따라 작성된다.
  \item 재무회계의 개념체계는 회계이론의 기본적 틀을 제공하는 것뿐만 아니라 실무에서 적용되는 회계실무의 기준을 보장한다.
  \item 재무보고에 있어서 중요성의 원칙은 거래유형에 관계없이 거래금액의 크기에 따라 판단한다.
  \item 재무제표의 질적특성 중 목적적합성과 신뢰성이 서로 상충될 경우 한쪽을 지나치게 강조하면 다른 쪽이 희생될 수 있다.
\end{choices}
\end{question}
\answer{5}
